\pdfcompresslevel=0

\catcode`\@=11
\def\@listi{\leftmargin\leftmargini
            \parsep=0pt
            \topsep=0pt
            \itemsep=0pt
						\partopsep=0pt}
\let\@listI\@listi
%\@listi
\def\@listii {\leftmargin\leftmarginii
              \labelwidth\leftmarginii
              \advance\labelwidth-\labelsep}
\def\@listiii{\leftmargin\leftmarginiii
              \labelwidth\leftmarginiii
              \advance\labelwidth-\labelsep}
\def\@listiv {\leftmargin\leftmarginiv
              \labelwidth\leftmarginiv
              \advance\labelwidth-\labelsep}
\def\@listv  {\leftmargin\leftmarginv
              \labelwidth\leftmarginv
              \advance\labelwidth-\labelsep}
\def\@listvi {\leftmargin\leftmarginvi
              \labelwidth\leftmarginvi
              \advance\labelwidth-\labelsep}
\catcode`\@=12

\baselineskip=0pt
\lineskiplimit=0pt
\lineskip=0pt
\parindent=0pt
\parskip=0pt
\topskip=0pt
\maxdepth=0pt
\pdfminorversion=7
\jot=4pt

% \pretolerance=-1
% \tolerance=9999
% \emergencystretch=8pt
\overfullrule=5mm
\showboxdepth=100
\showboxbreadth=1000

\def\capitalize#1#2{\uppercase{#1}#2}
\newbox\innerbox
\newbox\innerboxfirstline
\newbox\innerboxi
\newbox\partialpage
\newdimen\heightremaining
\newdimen\heightfirstline
\newdimen\innerwidth

\def\breakbox#1{%
	\global\def\breakboxcolor{#1}
	% TODO: Why do I need this? There is a problem when the page start with a breakbox. The whole line is not printed.
	% \hbox{}
	% First, save in \partialpage the content of the current page so far.
	% \vbox{\unvbox255} is used instead of \box255 because the height of the
	% \box255 is \textheight = \vsize (\pagegoal?)
	\output{\global\setbox\partialpage=\vbox{\unvbox255}}\eject
	\output{\myoutput}%
	% compute the width of the innerbox
	\innerwidth=\linewidth
	\advance\innerwidth by-2\fboxsep
	\advance\innerwidth by-2\fboxrule
	% The innerbox is a vbox that will contain the content of the theorem
	% It start a bgroup that will end at \endtheorem
	\setbox\innerbox=\vbox\bgroup%
		\hsize=\innerwidth
		\linewidth=\innerwidth
		% TODO: error when page breaking and the color was changed
}

\def\endbreakbox{%
	% It close the innerbox that was started at \theoremii
	\egroup
	% Config for the vsplit's
	\splitmaxdepth=0pt
	\splittopskip=0pt
	% Get the first line of the innerbox by splitting a copy
	\setbox\innerboxfirstline=\copy\innerbox
	\vfuzz=\maxdimen
	\setbox\innerboxfirstline=\vsplit\innerboxfirstline to0pt
	\vfuzz=1pt
	\setbox\innerboxfirstline=\vbox{\hsize=\innerwidth\unvbox\innerboxfirstline}%
	\heightfirstline=\ht\innerboxfirstline
	\advance\heightfirstline by\dp\innerboxfirstline
	% Compute the height remaining in the page
	\heightremaining=\textheight
	\advance\heightremaining by-\ht\partialpage
	\advance\heightremaining by-\dp\partialpage
	\advance\heightremaining by-2\fboxsep
	\advance\heightremaining by-2\fboxrule
	% Put the partial page back
	\box\partialpage
	% Check if the first line fits
	\ifdim\heightremaining>\heightfirstline
		% \message{First line fits}
		% If the firstline fits, then check if the whole box fits
		\ifdim\heightremaining<\ht\innerbox
			% \message{The whole box does not fit}
			% If The whole box does not fit, then split it
			\vbadness=10000
			\setbox\innerboxi=\vsplit\innerbox to\heightremaining
			\vbadness=1000
			% \vbox\unvbox\innerboxi instead of \box\innerboxi
			\setbox\innerboxi=\vbox{\unvbox\innerboxi}
			% Set the width of the innerbox so that the fcolorbox don't overfull
			\wd\innerboxi=\innerwidth
			\par\noindent\fcolorbox{\breakboxcolor}{\breakboxcolor!5}{\box\innerboxi}%
			\eject
		\else
			% \message{The whole box fits}
		\fi
	\else
		% If the firstline does not fit then eject to finish the current column
		% \message{First line does not fit}
		\eject
	\fi
	% If the innerbox is void after the splitting, then print it
	\ifvoid\innerbox\else
		% \message{The innerbox is not void}
		% Set the width of the innerbox so that the fcolorbox don't overflow
		\wd\innerbox=\innerwidth
		\smallskip
		\par\noindent\fcolorbox{\breakboxcolor}{\breakboxcolor!5}{\box\innerbox}\par
		\smallskip
	\fi
	\par
	\let\breakboxcolor=\relax
}

\newcount\thmnum
\thmnum=0
\def\newthm#1#2#3{
	% \theorem
	\expandafter\def\csname #3\endcsname{%
		\global\advance\thmnum by1
		% \message{\number\thmnum}%
		\def\tmp{\futurelet\next}%
		\expandafter\tmp\csname #3i\endcsname
	}

	% \theoremi
	\expandafter\def\csname #3i\endcsname{%
		\let\thmname=\relax
		\def\tmpi[####1]{\def\thmname{####1}\csname #3ii\endcsname}%
		\ifx\next[
			\expandafter\tmpi
		\else
			\csname #3ii\endcsname
		\fi
	}

	% \theoremii
	\expandafter\def\csname #3ii\endcsname{%
		\breakbox{#1}
		{\bf \capitalize#3 \number\thmnum.}~\ifx\thmname\relax\else(\thmname)\fi%
	}

	% \endtheorem
	\expandafter\def\csname end#3\endcsname{%
		\endbreakbox
	}
}


\edef\twocolumn{\detokenize{twocolumn}}
\edef\fourcolumn{\detokenize{fourcolumn}}
\edef\xjobname{\jobname}

\newcount\numcol
\ifx\xjobname\twocolumn
	\AtBeginDocument{
		\pdfpageheight=297mm
		\pdfpagewidth=210mm
	}
	\numcol=2
\else\ifx\xjobname\fourcolumn
	\numcol=4
	\AtBeginDocument{
		\pdfpageheight=297mm
		\pdfpagewidth=420mm
		\mag=707
	}
\else
	\numcol=1
	\ifx\documentclass\undefined
		\pdfpageheight=288mm
		\pdfpagewidth=100mm
	\else
		\AtBeginDocument{
			\pdfpageheight=290mm
			\pdfpagewidth=100mm
		}
	\fi
\fi\fi

\hoffset=-1in
\voffset=-1in


\countdef\pageno=0 \pageno=1
\newcount\curcol
\newcount\colbox
\newcount\firstbox
\firstbox=100
\curcol=0
\def\myoutput{%
	\global\advance\curcol by1
	{%
		\advance\curcol by\firstbox
		\global\setbox\curcol=\box255
		\global\wd\curcol=\textwidth
	}%
	\ifnum\curcol=\numcol
		\global\curcol=4
		\shipout\vbox to \pdfpageheight {%
			\vfil
			\hbox to \pdfpagewidth {%
				\colbox=0
				\hfil
				\loop
					\ifnum\colbox<\numcol
					\advance\colbox by1
					{%
						\hfil
						\advance\colbox by\firstbox
						\box\colbox
					}%
				\repeat
				\hfil\hfil
			}%
			\vfil
		}%
		\global\advance\pageno by1
		\global\curcol=0
	\fi
}

\def\xthm#1#2#3,{
	\if\relax\detokenize{#3}\relax\else
		\newthm{#1}{#2}{#3}
		\def\tmpxthm{\xthm{#1}{#2}}
		\expandafter\tmpxthm
	\fi
}

\ifx\documentclass\undefined\else
  \fboxsep=2mm
  \fboxrule=1.5pt
	\usepackage{xcolor}
	\definecolor{red}{RGB}{247,87,130}
	\definecolor{green}{RGB}{86,162,43}
	\definecolor{blue}{RGB}{126,134,244}
	\definecolor{bts}{RGB}{244,58,232}
	\definecolor{orange}{RGB}{235,106,44}
	\definecolor{gray}{RGB}{145,145,145}
	\xthm{green}{green}theorem,lemma,corollary,proposition,,
	\xthm{orange}{orange}definition,,
	\xthm{blue}{blue}question,remark,note,notation,claim,summary,acknowledgment,case,conclusion,observation,recall,affirmation,,
	\xthm{bts}{bts}example,application,problem,,
	\xthm{red}{red}exercise,,
	\xthm{gray}{gray}proof,solution,,

	\textheight=287mm
	\textwidth=97.5mm

\AtBeginDocument{
	\baselineskip=12pt
	\abovedisplayskip=1pt
	\belowdisplayskip=1pt
	\abovedisplayshortskip=2pt
	\belowdisplayshortskip=2pt
	\output{\myoutput}%
	\overfullrule=5mm
}
\fi

\def\coloneqq{\mathrel{\mathop:}=}
\def\mathclap#1{\hbox to 0pt{\hss$#1$\hss}}
\def\mathllap{}
